% ============================================================
% Chapter 3: Distributed Programming
% ============================================================
\section{分布式编程}

\subsection{从配置中启动 MQ}

\begin{enumerate}[leftmargin=1.5em]
    \item \lstinline|.erlang.cookie| 文件必须一致
    \item 必须一次性成功，把节点全部启动起来
\end{enumerate}

\subsection{Hidden Node}

\begin{figure}[H]
\centering
\includegraphics[width=0.85\textwidth]{image1}
\caption{Hidden Node}
\end{figure}

参考：\href{https://www.erlang.org/doc/reference_manual/distributed.html}{Erlang -- Distributed Erlang}

\subsection{远程连接 Erlang 节点}

\subsubsection{使用 remsh}

远程连接的最简方式。

\subsubsection{使用 net\_kernel:connect\_node}

\begin{lstlisting}[style=shellstyle]
[root@kvm_10_21_1_163 sbin]# erl -sname apple
(apple@kvm_10_21_1_163)1> net_kernel:connect_node('rabbit@kvm_10_20_1_60').
true
\end{lstlisting}

\begin{warnbox}[注意]
\lstinline|net_kernel:connect_node/1| 的返回值可能被忽略，参见 \href{https://blog.differentpla.net/blog/2013/11/07/net-kernel-connect-node-1-returns-ignored/}{net\_kernel:connect\_node/1 returns ignored}
\end{warnbox}

需要给节点起个名字（\lstinline|-sname|）才能连接。

\subsubsection{使用 net\_adm:ping}

注意是 \lstinline|-sname| 生效，不用 \lstinline|-name|，能够在远程节点生效。

\begin{figure}[H]
\centering
\includegraphics[width=0.85\textwidth]{image5}
\caption{net\_adm:ping 远程连接}
\end{figure}

\subsubsection{使用 rpc:call}

\begin{figure}[H]
\centering
\includegraphics[width=0.85\textwidth]{image6}
\caption{rpc:call 远程调用}
\end{figure}

\subsubsection{使用消息发送 !}

\begin{figure}[H]
\centering
\includegraphics[width=0.85\textwidth]{image7}
\caption{消息发送 ! 远程通信}
\end{figure}

\begin{figure}[H]
\centering
\includegraphics[width=0.85\textwidth]{image8}
\caption{消息发送示例}
\end{figure}

\subsubsection{使用 monitor}

\begin{figure}[H]
\centering
\includegraphics[width=0.8\textwidth]{image9}
\caption{monitor 用法 (1)}
\end{figure}

\begin{figure}[H]
\centering
\includegraphics[width=0.6\textwidth]{image10}
\caption{monitor 用法 (2)}
\end{figure}

\begin{figure}[H]
\centering
\includegraphics[width=0.7\textwidth]{image11}
\caption{monitor 用法 (3)}
\end{figure}

\begin{figure}[H]
\centering
\includegraphics[width=0.65\textwidth]{image12}
\caption{monitor 用法 (4)}
\end{figure}

\begin{figure}[H]
\centering
\includegraphics[width=0.65\textwidth]{image13}
\caption{monitor 用法 (5)}
\end{figure}

三种情况：
\begin{enumerate}[leftmargin=1.5em]
    \item Case 1: 节点不存在
    \item Case 2: 节点存在，服务不存在，建立连接
    \item Case 3: 远程节点和进程都存在
\end{enumerate}

\subsubsection{使用 Ctrl+G}

通过 Ctrl+G 可以连接到远程节点。

\subsubsection{使用 Observer 观察远程节点}

注意在本地先启动 observer，然后在 Nodes 中切换。

\textbf{步骤：}

\begin{enumerate}[leftmargin=1.5em]
    \item 新起一个节点
\end{enumerate}

\begin{figure}[H]
\centering
\includegraphics[width=0.85\textwidth]{image14}
\caption{新起节点}
\end{figure}

\begin{enumerate}[leftmargin=1.5em]
    \setcounter{enumi}{1}
    \item 启动 observer
\end{enumerate}

\begin{figure}[H]
\centering
\includegraphics[width=0.85\textwidth]{image15}
\caption{启动 observer}
\end{figure}

\begin{enumerate}[leftmargin=1.5em]
    \setcounter{enumi}{2}
    \item 连接其他 node
\end{enumerate}

\begin{tipbox}[提示]
在本机上即使不连接，能看到其他节点的时候也可以直接切换。在远程，必须先连接上，然后再切换。
\end{tipbox}

\begin{figure}[H]
\centering
\includegraphics[width=0.85\textwidth]{image16}
\caption{连接其他节点}
\end{figure}

或者用命令行：

\begin{figure}[H]
\centering
\includegraphics[width=0.85\textwidth]{image17}
\caption{命令行连接}
\end{figure}

\begin{enumerate}[leftmargin=1.5em]
    \setcounter{enumi}{3}
    \item 切换到其他节点
\end{enumerate}

\begin{figure}[H]
\centering
\includegraphics[width=0.85\textwidth]{image18}
\caption{切换到其他节点}
\end{figure}

\begin{enumerate}[leftmargin=1.5em]
    \setcounter{enumi}{4}
    \item 查看节点信息
\end{enumerate}

\begin{figure}[H]
\centering
\includegraphics[width=0.85\textwidth]{image19}
\caption{节点信息}
\end{figure}

\subsection{节点无法连接的调试}

\begin{itemize}[leftmargin=1.5em]
    \item 查看防火墙，如果有，需要关闭
    \item 检查 \lstinline|.erlang.cookie| 是否一致
    \item 检查 hosts 配置
\end{itemize}

\begin{figure}[H]
\centering
\includegraphics[width=0.85\textwidth]{image20}
\caption{节点连接调试 (1)}
\end{figure}

\begin{figure}[H]
\centering
\includegraphics[width=0.85\textwidth]{image21}
\caption{节点连接调试 (2)}
\end{figure}

\begin{figure}[H]
\centering
\includegraphics[width=0.85\textwidth]{image22}
\caption{节点连接调试 (3)}
\end{figure}

\subsection{把一个节点变成分布式节点}

为节点起名字即可：

\begin{lstlisting}
net_kernel:start([foo, shortnames]).
\end{lstlisting}

启动之前：

\begin{figure}[H]
\centering
\includegraphics[width=0.6\textwidth]{image30}
\caption{启动分布式之前}
\end{figure}

启动之后：

\begin{figure}[H]
\centering
\includegraphics[width=0.6\textwidth]{image31}
\caption{启动分布式之后}
\end{figure}

\begin{figure}[H]
\centering
\includegraphics[width=0.85\textwidth]{image32}
\caption{net\_kernel:start 结果}
\end{figure}

这是多出来的一部分：

\begin{figure}[H]
\centering
\includegraphics[width=0.85\textwidth]{image33}
\caption{新增部分 (1)}
\end{figure}

\begin{figure}[H]
\centering
\includegraphics[width=0.85\textwidth]{image34}
\caption{新增部分 (2)}
\end{figure}

源码位置：\lstinline|kernel-8.4.2/src/erl_distribution.erl|

\begin{figure}[H]
\centering
\includegraphics[width=0.85\textwidth]{image35}
\caption{erl\_distribution 源码 (1)}
\end{figure}

\begin{figure}[H]
\centering
\includegraphics[width=0.85\textwidth]{image36}
\caption{erl\_distribution 源码 (2)}
\end{figure}

\subsection{Erlang 节点间通信}

参见：\href{https://doc.weixin.qq.com/doc/w3_AIUAewaBACcKvp9d0QlSbafbqP8YO?scode=AJEAIQdfAAo4VWNB0wAIUAewaBACc}{erlang 节点间通信 -- Net Tick Time}

\subsection{找到 pid/reference/port 所在的 node}

参考：\href{https://www.erlang.org/doc/man/erlang.html#node-1}{Erlang -- erlang:node/1}
